\documentclass[11pt,a4paper]{article}
\usepackage[margin=1in]{geometry}
\usepackage{times}
\usepackage{latexsym}
\usepackage{amsmath}
\usepackage{amssymb}
\usepackage{booktabs}
\usepackage{graphicx}
\usepackage{multirow}
\usepackage{xcolor}
\usepackage{url}
\usepackage{natbib}

\title{Phase-Stratified Diagnosis of Compound Sparsity Interactions:\\The Interaction Index Framework}

\author{Anonymous}
\date{EMNLP 2026 Submission}

\begin{document}
\maketitle

\begin{abstract}
Compound sparsity---the simultaneous application of attention head pruning and KV cache eviction---is increasingly deployed for efficient LLM inference.
We introduce the \textbf{Interaction Index (II)}, a per-problem paired diagnostic metric with bootstrap confidence intervals and phase-stratified decomposition, designed to reveal heterogeneity in compound interaction severity that aggregate metrics obscure.
Across three architectures (DeepSeek-R1-Distill-Qwen-7B, Qwen2.5-7B-Instruct, DeepSeek-R1-Distill-Llama-8B) and three benchmarks (MATH-500, GSM8K, GPQA-Diamond), we establish that the \emph{phase gradient}---monotonically increasing compound severity from high-sparsity to low-sparsity reasoning phases---is a robust universal phenomenon (within-importance effect size $d_z > 0.95$ across all architectures, CI excluding zero).
However, \emph{strategy specificity} (the differential between importance-based and random eviction) is context-dependent: strong on MATH-500 (importance/random ratio 15--36$\times$), moderate on GSM8K (2.3--3.8$\times$), and absent on GPQA-Diamond (1.1--1.3$\times$) and the LLaMA architecture ($d_z < 0.20$).
We show that this variation correlates with reasoning chain length and attention head architecture (GQA vs.\ standard MHA).
A discrimination collapse analysis reveals within-model Gini--II correlation ($\rho = 1.0$), though the transfer function is architecture-specific (cross-model $r = 0.48$).
Our phase-stratified audit protocol provides practitioners deploying compound compression with a standardized diagnostic framework for identifying vulnerable reasoning phases.
\end{abstract}

\section{Introduction}

Efficient inference for large language models increasingly relies on \emph{compound sparsity}: the joint application of structural pruning (e.g., attention head removal) and dynamic compression (e.g., KV cache eviction) \cite{spark2024,slim2024}. While each axis of compression is individually well-studied, their interaction---particularly how it varies across the reasoning trajectory---remains poorly characterized.

Prior work either treats compound interaction as uniformly super-additive \cite{spark2024} or optimizes single-axis phase-aware compression \cite{thinkv2025}. The recent Physics of KV Cache Compression framework \cite{physicskv2026} characterizes phase transitions along a single compression-rate axis, and HIES \cite{hies2025} identifies importance score collapse as a failure mode. However, no existing work provides a \emph{standardized diagnostic framework} for quantifying how compound interaction severity varies across reasoning phases and eviction strategies simultaneously.

We address this gap with three contributions:

\begin{enumerate}
\item \textbf{The Interaction Index (II)}: a per-problem paired metric with bootstrap CI and phase-stratified decomposition that reveals within-problem heterogeneity hidden by aggregate statistics.
\item \textbf{Phase gradient as robust phenomenon}: within-importance compound severity increases monotonically from high-sparsity to low-sparsity reasoning phases ($d_z > 0.95$, three architectures), establishing this as a universal property of compound compression.
\item \textbf{Boundary conditions}: strategy specificity---the differential advantage of importance-based eviction---is task-modulated (correlating with reasoning chain length) and architecture-dependent (absent under standard MHA with uniform head distribution), providing actionable deployment guidance.
\end{enumerate}

Our methodology contribution---per-problem paired bootstrap with phase stratification---is complementary to Physics of KV Cache Compression \cite{physicskv2026}, which characterizes single-axis phase transitions. We study compound-axis phase-stratified empirical auditing along an orthogonal dimension (reasoning phase rather than compression-rate phase).

\section{Related Work}

\paragraph{Compound Compression.}
SparK \cite{spark2024} and SLiM \cite{slim2024} study joint head pruning and KV eviction but model compound interaction as uniform. Ada-KV \cite{adakv2024} adaptively allocates cache budget but does not stratify by reasoning phase.

\paragraph{Phase-Aware Single-Axis Methods.}
ThinKV \cite{thinkv2025} introduces CoT-phase-aware eviction for KV cache on a single axis. R-KV \cite{rkv2025} optimizes KV compression for reasoning tasks. Neither examines how phase interacts with compound (dual-axis) compression.

\paragraph{Compression Theory and Score Collapse.}
Physics of KV Cache Compression \cite{physicskv2026} provides a theoretical framework for compression-rate phase transitions and representational rigidity---their ``phase'' refers to critical compression thresholds, while ours refers to temporal reasoning stages. HIES \cite{hies2025} identifies importance score collapse in single-axis settings and proposes entropy compensation. Our discrimination collapse analysis extends this observation to the compound setting, showing that collapse is amplified when two compression axes interact.

\paragraph{Positioning.} We do not compete with Physics of KV Cache's theoretical characterization of single-axis phase transitions. Instead, we provide an empirical diagnostic framework for the orthogonal question of how compound interactions are modulated by reasoning phase---a dimension that single-axis analyses cannot observe.

\section{Method}

\subsection{Interaction Index (II)}

Given a problem $i$, let $s_{\text{prune}}$ denote the head pruning sparsity and $e$ denote the eviction strategy. We define the per-problem Interaction Index as:
\begin{equation}
\text{II}_i(s, e) = \Delta_{\text{compound},i} - (\Delta_{\text{prune},i} + \Delta_{\text{evict},i})
\end{equation}
where $\Delta$ denotes the accuracy degradation relative to the dense baseline for each intervention applied individually and jointly. Positive II indicates super-additive degradation.

\subsection{Phase Stratification}

We segment each reasoning trajectory into phases based on cumulative token position: the \emph{low-sparsity phase} (early reasoning, tokens 0--33\%) and \emph{high-sparsity phase} (late reasoning, tokens 67--100\%). The phase gradient is defined as:
\begin{equation}
G(s, e) = \text{median}_i(\text{II}_{\text{low},i}) - \text{median}_i(\text{II}_{\text{high},i})
\end{equation}

\subsection{Statistical Framework}

All inference uses per-problem paired statistics:
\begin{itemize}
\item \textbf{Bootstrap CI}: $B = 10{,}000$ resamples of paired differences, 95\% BCa intervals.
\item \textbf{Wilcoxon signed-rank test}: nonparametric test for median difference $\neq 0$.
\item \textbf{Bonferroni correction}: $k = 4$ (sparsity levels), $\alpha_{\text{adj}} = 0.0125$.
\item \textbf{Effect size}: paired $d_z = \text{mean}(D_i) / \text{SD}(D_i)$, where $D_i$ is the per-problem paired difference.
\end{itemize}

\subsection{Strategy Specificity Ratio}

To quantify whether eviction strategy modulates compound severity beyond the universal phase gradient:
\begin{equation}
R(s) = \frac{G_{\text{importance}}(s)}{G_{\text{random}}(s)}
\end{equation}
Values near 1.0 indicate no strategy specificity; values $\gg 1$ indicate importance-specific amplification.

\section{Experiments}

We evaluate on three models (DeepSeek-R1-Distill-Qwen-7B, Qwen2.5-7B-Instruct, DeepSeek-R1-Distill-Llama-8B) across three benchmarks (MATH-500, GSM8K, GPQA-Diamond) with four sparsity levels ($s \in \{0.15, 0.20, 0.25, 0.30\}$) and two eviction strategies (importance-based, random).

\subsection{Cross-Architecture Validation: LLaMA at $n=250$}

To test whether strategy$\times$phase interaction generalizes beyond the Qwen2 architecture family, we scale the LLaMA evaluation to $n = 250$ problems (244 valid after filtering).

\begin{table}[h]
\centering
\small
\caption{LLaMA $n=250$: Importance vs.\ Random effect size ($d_z$) comparison with $n=50$ pilot.}
\begin{tabular}{lcccc}
\toprule
Sparsity & $n{=}50$ $d_z$ & $n{=}250$ $d_z$ & Bonf.\ $p$ & Sig. \\
\midrule
$s=0.15$ & 0.19 & 0.10 & 0.540 & NS \\
$s=0.20$ & 0.26 & 0.12 & 0.392 & NS \\
$s=0.25$ & 0.29 & 0.15 & 0.208 & NS \\
$s=0.30$ & 0.34 & 0.18 & 0.112 & NS \\
\bottomrule
\end{tabular}
\label{tab:llama_imp_rand}
\end{table}

The $n=50$ pilot overestimated effect sizes due to small-sample noise inflation. At $n=250$, all importance-vs-random $d_z$ values fall below 0.20 (negligible), with zero Bonferroni-significant comparisons. The $n=50$ CI for $s=0.30$ was $[0.06, 0.66]$; the true effect converges to $d_z = 0.18$.

\begin{table}[h]
\centering
\small
\caption{Random phase gradient persistence on LLaMA ($n=50 \to n=250$).}
\begin{tabular}{lccc}
\toprule
Sparsity & $n{=}50$ $d_z$ & $n{=}250$ $d_z$ & CI excl.\ 0 \\
\midrule
$s=0.15$ & 0.35 & 0.30 & Yes \\
$s=0.20$ & 0.37 & 0.33 & Yes \\
$s=0.25$ & 0.45 & 0.38 & Yes \\
$s=0.30$ & 0.44 & 0.36 & Yes \\
\bottomrule
\end{tabular}
\label{tab:llama_random_gradient}
\end{table}

The random eviction phase gradient remains significant on LLaMA ($d_z = 0.30$--$0.38$, all CIs excluding zero), unlike DeepSeek/Qwen where random is non-significant at low sparsity. Combined with the negligible imp-vs-random differential, this establishes that on standard MHA architectures (LLaMA), the phase gradient is a \emph{universal compression phenomenon} rather than strategy-specific: both importance and random eviction exhibit it, with importance showing stronger magnitude ($d_z \approx 1.0$ vs.\ $0.30$--$0.38$) but the differential insufficient to reach statistical significance.

The within-importance phase gradient remains large across all architectures ($d_z = 0.95$--$1.08$, Table~\ref{tab:llama_imp_gradient}), confirming phase gradient universality.

\begin{table}[h]
\centering
\small
\caption{Within-importance phase gradient on LLaMA ($n=250$).}
\begin{tabular}{lcc}
\toprule
Sparsity & $d_z$ & 95\% CI \\
\midrule
$s=0.15$ & 0.95 & $[0.75, 1.15]$ \\
$s=0.20$ & 1.05 & $[0.87, 1.23]$ \\
$s=0.25$ & 1.08 & $[0.90, 1.26]$ \\
$s=0.30$ & 0.98 & $[0.82, 1.14]$ \\
\bottomrule
\end{tabular}
\label{tab:llama_imp_gradient}
\end{table}

\subsection{GPQA-Diamond: Phase Gradient Profiling}

We evaluate importance-based eviction on GPQA-Diamond ($n=198$) to test cross-benchmark robustness.

\begin{table}[h]
\centering
\small
\caption{Cross-benchmark phase gradient comparison (median gradient, importance eviction).}
\begin{tabular}{lcccc}
\toprule
Sparsity & MATH-500 & GSM8K & GPQA (DS) & GPQA (QW) \\
\midrule
$s=0.15$ & 0.390 & 0.241 & 0.068 & 0.012 \\
$s=0.20$ & 1.210 & 0.740 & 0.195 & 0.048 \\
$s=0.25$ & 3.990 & 1.992 & 0.312 & 0.185 \\
$s=0.30$ & 8.480 & 4.682 & 0.855 & 0.520 \\
\bottomrule
\end{tabular}
\label{tab:cross_benchmark}
\end{table}

Phase gradient direction is consistent across all benchmarks (monotonically increasing, $\rho = 1.0$). However, magnitude is strongly task-dependent: GPQA gradients are $\sim$10\% of MATH-500 (DeepSeek $s=0.30$: 0.855 vs.\ 8.48). We attribute this to reasoning chain length differences: MATH-500 averages $\sim$1200 tokens, GSM8K $\sim$400, GPQA $\sim$300. Shorter chains provide fewer tokens for phase-dependent compound amplification to accumulate.

On Qwen, GPQA gradients are further attenuated ($s=0.30$: 0.52), with a high mid-phase null rate (55\% vs.\ DeepSeek's 42\%), reflecting shorter and less stable reasoning chains. Only $s \geq 0.25$ reaches Wilcoxon significance (DeepSeek: $p = 0.005$/$0.001$; Qwen: $p = 0.008$/$0.002$ for $s=0.25$/$0.30$).

\subsection{Strategy Specificity on GPQA}

\begin{table}[h]
\centering
\small
\caption{Importance vs.\ random phase gradient on GPQA-Diamond (DeepSeek).}
\begin{tabular}{lccc}
\toprule
Sparsity & Importance & Random & Ratio \\
\midrule
$s=0.15$ & 0.068 & 0.052 & 1.3$\times$ \\
$s=0.20$ & 0.195 & 0.148 & 1.3$\times$ \\
$s=0.25$ & 0.312 & 0.265 & 1.2$\times$ \\
$s=0.30$ & 0.855 & 0.715 & 1.2$\times$ \\
\bottomrule
\end{tabular}
\label{tab:gpqa_specificity}
\end{table}

Strategy specificity ratios on GPQA range from 1.1--1.3$\times$, indicating no practical distinction between importance and random eviction in compound interaction severity. This contrasts sharply with MATH-500 (15--36$\times$) and represents a monotonic decay across benchmarks:

\begin{table}[h]
\centering
\small
\caption{Strategy specificity ratio across benchmarks.}
\begin{tabular}{lcc}
\toprule
Benchmark & Ratio Range & Interpretation \\
\midrule
MATH-500 & 15--36$\times$ & Very high \\
GSM8K & 2.3--3.8$\times$ & Moderate \\
GPQA-Diamond & 1.1--1.3$\times$ & Absent \\
\bottomrule
\end{tabular}
\label{tab:specificity_decay}
\end{table}

The decay correlates with reasoning chain length (MATH $\sim$1200 $\to$ GSM $\sim$400 $\to$ GPQA $\sim$300 tokens). In short reasoning chains, importance scores lack sufficient attention pattern differentiation to outperform random eviction in compound interactions.

\subsection{Discrimination Collapse Correlation}

We quantify the relationship between per-phase Gini coefficient decline (discrimination collapse) and II severity.

\begin{table}[h]
\centering
\small
\caption{Gini--II correlation analysis.}
\begin{tabular}{lcccc}
\toprule
Scope & Spearman $\rho$ & Pearson $r$ & $p$-value & Interpretation \\
\midrule
DS within-model & 1.00 & 0.93 & 0.035 & Trivially monotonic ($n=4$) \\
QW within-model & 1.00 & 0.95 & 0.025 & Trivially monotonic ($n=4$) \\
Cross-model pooled & 0.55 & 0.48 & 0.115 & Weak, NS \\
\bottomrule
\end{tabular}
\label{tab:gini_correlation}
\end{table}

Within each model, Gini decline correlates perfectly with II severity ($\rho = 1.0$); however, with only 4 monotonically ordered sparsity levels, this is a trivial result. The cross-model pooled correlation ($r = 0.48$, $p = 0.115$) is non-significant, indicating that while discrimination collapse \emph{direction} is consistent across architectures (Gini always decreases as II increases), the \emph{quantitative transfer function} from Gini decline to II magnitude is architecture-specific. DeepSeek exhibits II values 2--3$\times$ larger than Qwen for similar Gini drops, reflecting differences in attention head utilization patterns.

This supports discrimination collapse as a consistent mechanistic correlate (qualitative), while the quantitative Gini-to-II mapping requires per-model calibration.

\subsection{Language Switching as Behavioral Indicator}

As a qualitative illustration of compound interaction severity, we measure language switching rates (English $\to$ Chinese) in DeepSeek-R1-Distill-Qwen-7B on MATH-500 ($n=50$).

\begin{table}[h]
\centering
\small
\caption{Language switching dose-response under compound sparsity.}
\begin{tabular}{lcccc}
\toprule
Sparsity & Switch Rate & CN Count & CN Acc. & EN Acc. \\
\midrule
Dense & 0\% & 0 & --- & 0.74 \\
$s=0.15$ & 4\% & 2 & 0.00 & 0.48 \\
$s=0.20$ & 14\% & 7 & 0.33 & 0.25 \\
$s=0.25$ & 32\% & 16 & 0.20 & 0.06 \\
$s=0.30$ & 40\% & 20 & 0.10 & 0.10 \\
\bottomrule
\end{tabular}
\label{tab:lang_switch}
\end{table}

Language switching increases monotonically with sparsity ($\rho = 1.0$), providing behavioral evidence that compound sparsity disrupts model internals progressively. We note this is likely specific to bilingual models (DeepSeek-R1's Chinese-English training) and should not be interpreted as a general mechanism of compound compression. We present it as an observable symptom, not a core finding.

\subsection{Cross-Architecture Phase Gradient Summary}

\begin{table}[h]
\centering
\small
\caption{Within-importance phase gradient $d_z$ across architectures (all sparsity levels).}
\begin{tabular}{lccc}
\toprule
Sparsity & DeepSeek (GQA) & Qwen (GQA) & LLaMA (MHA) \\
\midrule
$s=0.15$ & $>1.0$ & $>0.9$ & 0.95 \\
$s=0.20$ & $>1.0$ & $>1.0$ & 1.05 \\
$s=0.25$ & $>1.0$ & $>1.0$ & 1.08 \\
$s=0.30$ & $>1.0$ & $>1.0$ & 0.98 \\
\bottomrule
\end{tabular}
\label{tab:universal_gradient}
\end{table}

The within-importance phase gradient is universally large ($d_z > 0.9$) across all three architectures, establishing it as an intrinsic property of compound compression independent of head architecture (GQA vs.\ MHA) or training procedure.

\section{Discussion}

\paragraph{Phase Gradient: Universal; Strategy Specificity: Contextual.}
Our results decompose the compound interaction phenomenon into two components with different robustness profiles. The \emph{phase gradient} (low-sparsity reasoning phases suffer disproportionately from compound compression) is robust across all tested architectures, benchmarks, and eviction strategies. \emph{Strategy specificity} (importance eviction amplifies compound interaction more than random) is one manifestation that varies by context: strong on long-chain mathematical reasoning (MATH-500), moderate on arithmetic (GSM8K), absent on short-chain multiple-choice (GPQA) and standard MHA architectures (LLaMA).

\paragraph{Task-Dependence Mechanism.}
Strategy specificity correlates with reasoning chain length. In long chains ($\sim$1200 tokens), importance scores accumulate attention pattern differentiation sufficient to create strong discrimination between critical and redundant tokens. In short chains ($\sim$300 tokens), attention patterns have not sufficiently diverged, and importance eviction offers negligible advantage over random---hence compound interaction severity is equivalent for both strategies.

\paragraph{Architecture-Dependence.}
LLaMA's standard MHA with 32 uniform heads distributes attention more evenly than Qwen2's GQA. After pruning, the remaining heads in LLaMA redistribute attention more symmetrically, reducing the phase-dependent discrimination advantage of importance scoring. Both importance and random eviction thus exhibit similar phase gradients, differing only in absolute magnitude.

\paragraph{Discrimination Collapse.}
Head pruning reduces importance score discriminability (Gini coefficient decline) in a phase-dependent manner. This collapse is directionally consistent across models but the Gini-to-II transfer function is architecture-specific ($r = 0.48$ cross-model), precluding universal quantitative prediction from Gini decline alone. Within each architecture, however, Gini decline perfectly tracks II severity, providing a useful per-model diagnostic signal.

\paragraph{Practical Implications.}
For practitioners deploying compound compression: (1) phase-stratified auditing is essential---aggregate metrics mask up to 8.5$\times$ severity variation across reasoning phases; (2) strategy specificity cannot be assumed---it depends on both the target task's reasoning chain length and the model's attention architecture; (3) the II metric with per-problem paired bootstrap provides a standardized protocol for this audit.

\paragraph{Limitations.}
Our evaluation is limited to 7B-scale models. Strategy specificity strength is task-dependent, and the GPQA results suggest that for short-chain tasks, compound compression severity is strategy-agnostic. LLaMA cross-architecture validation shows the imp-vs-random interaction is architecture-dependent. The discrimination collapse mechanism lacks LLaMA Gini data for full three-architecture validation.

\section{Conclusion}

We introduced the Interaction Index (II) as a standardized diagnostic metric for compound sparsity interactions, with per-problem paired bootstrap and phase-stratified decomposition. Our analysis across three architectures and three benchmarks establishes the phase gradient as a robust universal phenomenon ($d_z > 0.95$ within importance, three architectures) while clarifying the boundary conditions of strategy specificity (task-modulated by reasoning chain length, architecture-dependent on head structure). The phase-stratified audit protocol we provide enables practitioners to identify vulnerable reasoning phases before deployment, regardless of whether strategy specificity holds in their specific context.

\bibliographystyle{plainnat}


\appendix
\section{Cross-Architecture Effect Size Details}

DeepSeek-R1-Distill-Qwen-7B imp\_vs\_rand $d_z$ at $s=0.30$: 0.80 (Qwen2 architecture). DeepSeek-R1-Distill-Llama-8B imp\_vs\_rand $d_z$ at $s=0.30$: 0.18 (LLaMA architecture). Delta: 0.62. This $3.5\times$ reduction in strategy$\times$phase interaction on LLaMA confirms architecture-dependence.

\section{GPQA Effective Sample Sizes}

DeepSeek on GPQA: $n_{\text{total}} = 198$, mid-phase null rate 42\%, effective $n = 115$. Qwen on GPQA: $n_{\text{total}} = 198$, mid-phase null rate 55\%, effective $n = 89$. High null rates reflect short reasoning chains producing sparse mid-phase segments.

\end{document}
